% TEMPLATE-MILC-v3.tex - plantilla maestra UNICA de las guias MILC v3.
%
% La usan las tres familias de guias, cada una con su driver:
%   · Grados 6-11  -> scripts/build-guias-g11.py
%   · Bebras       -> scripts/build-guias-bebras.py
%   · Semillero    -> scripts/build-guias-semillero.py
%
% Antes habia tres copias casi identicas (~400 lineas iguales, 6 distintas):
% cada mejora habia que aplicarla tres veces, y en la practica no se hacia
% (el motor de recursos vivia solo en dos de ellas). Lo que varia entre
% familias esta parametrizado en 6 placeholders que cada driver llena
% (se nombran aqui SIN su sintaxis real: el builder reemplaza por texto plano,
% tambien dentro de los comentarios, y un valor multilinea rompe el preambulo):
%
%   HEADER_IZQ        encabezado izquierdo de pagina
%   HEADER_DER        encabezado derecho de pagina
%   PORTADA_ETIQUETA  rotulo sobre el numero grande (GRADO/MOMENTO/MODULO)
%   PORTADA_SERIE     linea de serie de la portada
%   PORTADA_ID        identificador de la guia en la portada
%   PORTADA_CIRCULO   el circulito de grado (°), o vacio si no aplica
%
% El resto de claves las documenta content/guias/_SCHEMA.md. El script valida
% que no queden placeholders sin reemplazar antes de escribir el archivo final.

\documentclass[11pt,letterpaper]{article}
\usepackage[margin=1.75cm]{geometry}
\usepackage[table]{xcolor}
\usepackage{fontspec}
\usepackage[spanish]{babel}
\usepackage{tikz}
% Librerías TikZ para los diagramas de las guías (bloque `recursos:`):
%  · babel  → babel en español vuelve activos caracteres como «>», y TikZ los usa
%             en sus opciones (>=stealth, ->). Sin esta librería, cualquier
%             diagrama con flechas revienta con «\language@active@arg> has an
%             extra }». Debe ir ANTES de que se dibuje nada.
%  · positioning        → sintaxis «below left=2pt and 3pt».
%  · decorations.pathreplacing → llaves ({brace}) para señalar rangos.
\usetikzlibrary{babel,positioning,decorations.pathreplacing}
\usepackage{graphicx}
% Circuitos: el trabajo de electrónica se hace hoy con micro:bit y sus sensores
% integrados (MakeCode, sin cableado), así que no se necesitan esquemáticos.
% Si algún día entran montajes con protoboard, basta descomentar esta línea:
% los diagramas .tex/.tikz declarados en `recursos:` ya se incrustan solos.
% \usepackage{circuitikz}
\usepackage[most]{tcolorbox}
\usepackage{tabularx}
\usepackage{array}
\usepackage{fancyhdr}
\usepackage{titlesec}
\usepackage{ragged2e}
\usepackage{enumitem}
\usepackage{microtype}

% Helvetica Neue no trae la flecha U+2192, asi que cada «→» escrito literalmente
% en el YAML se compilaba a NADA, en silencio (462 flechas en 61 guias: rutas del
% tipo «paleta → tipografia → lockup» salian rotas). Mapeamos el caracter a su
% version matematica, que si tiene glifo. Asi el autor puede seguir escribiendo
% «→» comodamente en el YAML.
\usepackage{newunicodechar}
\newunicodechar{→}{\ensuremath{\rightarrow}}

\setmainfont{Helvetica Neue}
\setsansfont{Helvetica Neue}
\setlength{\parindent}{0pt}
\setlength{\parskip}{6pt}
\hyphenpenalty=200
\exhyphenpenalty=200
\pretolerance=400
\tolerance=2000
\setlength{\emergencystretch}{1em}
\renewcommand{\arraystretch}{1.25}
\setlist{leftmargin=5mm,itemsep=1mm,topsep=1mm}
\newcolumntype{Y}{>{\RaggedRight\arraybackslash}X}

\definecolor{milcMagenta}{HTML}{E6007E}
\definecolor{milcVino}{HTML}{5A0038}
\definecolor{milcTurquesa}{HTML}{0093A5}
\definecolor{milcVerde}{HTML}{58A923}
\definecolor{milcMostaza}{HTML}{E5B400}
\definecolor{milcOcre}{HTML}{B86600}
\definecolor{milcNegro}{HTML}{241816}
\definecolor{milcGris}{HTML}{F7F7F4}
\definecolor{milcLinea}{HTML}{E8E2DD}
\definecolor{milcGradePrimary}{HTML}{FF2D87}
\definecolor{milcGradeDark}{HTML}{9F1B56}
\definecolor{milcGradeSoft}{HTML}{FFE0EE}
\definecolor{milcGradeAccent}{HTML}{CFE7FF}
\definecolor{milcGradeText}{HTML}{FFFFFF}
\color{milcNegro}

\pagestyle{fancy}
\fancyhf{}
\lhead{\small Tecnología e Informática · Grado 10}
\rhead{\small Guía 12 · MILC v3}
\cfoot{\small\thepage}
\renewcommand{\headrulewidth}{0pt}

\newtcolorbox{titlebox}[1]{
  enhanced,colback=#1,colframe=#1,arc=7mm,boxrule=0pt,
  left=7mm,right=7mm,top=3mm,bottom=3mm,
  before skip=8pt,after skip=8pt,
  fontupper=\sffamily\bfseries\Large\color{white}
}

\newtcolorbox{softbox}[3]{
  enhanced,colback=#2,colframe=#1,borderline west={4pt}{0pt}{#1},
  arc=2mm,boxrule=.4pt,left=4mm,right=4mm,top=3mm,bottom=3mm,
  before skip=6pt,after skip=8pt,title={#3},coltitle=white,
  colbacktitle=#1,fonttitle=\sffamily\bfseries,fontupper=\RaggedRight,
  attach boxed title to top left={xshift=3mm,yshift=-2mm},
  boxed title style={arc=2mm, boxrule=0pt}
}

\newtcolorbox{drawbox}[2]{
  enhanced,colback=white,colframe=#1,arc=4mm,boxrule=1pt,
  left=5mm,right=5mm,top=4mm,bottom=4mm,before skip=8pt,after skip=8pt,
  title={#2},coltitle=white,colbacktitle=#1,fonttitle=\sffamily\bfseries,
  fontupper=\RaggedRight,
  attach boxed title to top left={xshift=4mm,yshift=-2mm},
  boxed title style={arc=3mm, boxrule=0pt}
}

\newtcolorbox{infoband}[1]{
  enhanced,colback=#1!8,colframe=#1!50,arc=2mm,boxrule=.5pt,
  left=4mm,right=4mm,top=2mm,bottom=2mm,before skip=4pt,after skip=4pt,
  fontupper=\small\RaggedRight
}

% Puente narrativo entre secciones (contrato editorial Conéctate).
% Da continuidad: "Acabas de X. Ahora vas a Y porque Z."
% Visualmente: banda izquierda en color del grado, texto en itálica pequeña.
\newtcolorbox{puentebox}{
  enhanced,colback=milcGradeSoft,colframe=milcGradeDark,
  borderline west={3pt}{0pt}{milcGradeDark},
  arc=0mm,boxrule=0pt,left=5mm,right=4mm,top=2.5mm,bottom=2.5mm,
  before skip=4pt,after skip=8pt,
  fontupper=\itshape\small\color{milcNegro}\RaggedRight
}
% La flecha va en modo matematico a proposito: Helvetica Neue NO tiene el glifo
% U+2192, asi que \textrightarrow se compilaba a NADA («Missing character: There
% is no → in font Helvetica Neue») y el puente salia sin flecha. En modo
% matematico la flecha viene de la fuente matematica y si se dibuja.
\newcommand{\puente}[1]{%
  \begin{puentebox}%
    \textcolor{milcGradeDark}{$\rightarrow$}\hspace{2mm}#1%
  \end{puentebox}%
}

\newcommand{\linead}[1]{\noindent\color{milcLinea}\rule{#1}{0.5pt}\color{milcNegro}}
\newcommand{\respuesta}[1]{\par\vspace{2mm}\foreach \n in {1,...,#1}{\noindent\color{milcLinea}\rule{\linewidth}{0.5pt}\par\vspace{5mm}}\color{milcNegro}}
\newcommand{\checkbox}{\textcolor{milcGradeDark}{\fbox{\phantom{x}}}\hspace{5pt}}

\newcommand{\phasecard}[4]{
\begin{tcolorbox}[enhanced,colback=#3,colframe=#2,arc=3mm,boxrule=.5pt,height=3.05cm,valign=center,halign=center,left=1.5mm,right=1.5mm,top=2mm,bottom=2mm]
{\sffamily\bfseries\fontsize{22}{24}\selectfont\textcolor{#2}{#1}}\par
{\sffamily\bfseries\small #4}
\end{tcolorbox}
}

% ── Recursos visuales de la guía ──────────────────────────────────────
% Declarados en el bloque `recursos:` del YAML y emitidos por el builder.
% \guiaFigura   → imagen rasterizada o PDF (foto, carta celeste, montaje).
% \guiaDiagrama → fuente TikZ/circuitikz (.tex) incrustada como vector:
%                 es la vía para circuitos y esquemáticos editables.
% #1 = ruta relativa al .tex · #2 = pie (puede ir vacío)
\newcommand{\guiaFigura}[2]{%
\begin{center}
\includegraphics[width=0.88\linewidth,height=0.38\textheight,keepaspectratio]{#1}\\[1.5mm]
{\footnotesize\itshape\textcolor{milcNegro!75}{#2}}
\end{center}
}

\newcommand{\guiaDiagrama}[2]{%
\begin{center}
\input{#1}\\[1.5mm]
{\footnotesize\itshape\textcolor{milcNegro!75}{#2}}
\end{center}
}

\begin{document}
\thispagestyle{empty}
\begin{tikzpicture}[remember picture,overlay]
  \fill[white] (current page.south west) rectangle (current page.north east);
  \path[fill=milcGradePrimary,rounded corners=16mm]
    ([xshift=.55cm,yshift=-.55cm]current page.north west) rectangle
    ([xshift=-.55cm,yshift=3.65cm]current page.south east);

  \node[anchor=north west,text=milcGradeText] at ([xshift=2.0cm,yshift=-1.85cm]current page.north west)
    {\sffamily\bfseries\fontsize{14}{16}\selectfont GRADO};
  \node[anchor=north west,text=milcGradeText,opacity=.82] at ([xshift=2.0cm,yshift=-2.75cm]current page.north west)
    {\sffamily\bfseries\fontsize{9}{11}\selectfont SERIE GUÍAS · TECNOLOGÍA E INFORMÁTICA};

  \begin{scope}[shift={([xshift=-3.05cm,yshift=-2.55cm]current page.north east)}]
    \fill[white,opacity=.80] (-.62,-.52) rectangle (.62,.52);
    \draw[milcGradeText,line width=1pt] (-.62,-.52) rectangle (.62,.52);
    \fill[milcGradeDark] (-.42,-.38) rectangle (-.22,.06);
    \fill[milcGradeAccent] (-.10,-.38) rectangle (.10,.24);
    \fill[milcGradePrimary!60!black] (.22,-.38) rectangle (.42,.38);
  \end{scope}

  \node[anchor=west,text=milcGradeText] at ([xshift=1.9cm,yshift=-12.85cm]current page.north west)
    {\sffamily\bfseries\fontsize{124}{124}\selectfont 10};
  % El circulito de grado (°) solo aplica a las guias de grado («11°»); en
  % Bebras y Semillero el numero grande es un momento/modulo, no un grado.
  \draw[milcGradeText,line width=1.7pt]
    ([xshift=4.52cm,yshift=-11.68cm]current page.north west) circle (.13cm);

  \node[anchor=west,text=milcGradeText,text width=8.7cm,align=left] at ([xshift=10.35cm,yshift=-11.6cm]current page.north west)
    {
     {\sffamily\bfseries\fontsize{19}{22}\selectfont Estructura del\\informe técnico ---\\intro, desarrollo, cierre}\\[4mm]
     {\sffamily\bfseries\fontsize{23}{26}\selectfont Guía 12-10-TIC}\\[2mm]
     {\sffamily\fontsize{13}{16}\selectfont Período 2 · Informes técnicos y comunicación profesional con IA}\\[5mm]
     {\sffamily\bfseries\fontsize{11}{13}\selectfont Institución Educativa Sor María Juliana}\\[1mm]
     {\sffamily\fontsize{10}{12}\selectfont Autor: Dr. Álvaro Cárdenas Orozco}
    };

  \path[fill=milcGradeSoft,rounded corners=7mm]
    ([xshift=1.35cm,yshift=.85cm]current page.south west) rectangle
    ([xshift=-1.35cm,yshift=4.25cm]current page.south east);
  \draw[milcGradeDark,line width=.65pt,rounded corners=7mm]
    ([xshift=1.35cm,yshift=.85cm]current page.south west) rectangle
    ([xshift=-1.35cm,yshift=4.25cm]current page.south east);
  \node[anchor=center,text=milcNegro,text width=17.0cm,align=center] at ([yshift=2.55cm]current page.south)
    {\sffamily\fontsize{10}{12}\selectfont Metodología MILC · Apertura ancestral · Pensamiento computacional · Triángulo Dussel-Estoicismo-Floridi\\
    \textcolor{milcGradeDark}{\bfseries Producto:} Plantilla personal reutilizable de informe técnico con las 3 secciones obligatorias estructuradas (introducción con contexto, objetivo y alcance; desarrollo con datos, argumentos y análisis; conclusiones con hallazgos y recomendaciones accionables) + 1 informe técnico aplicado a un caso real del colegio (propuesta de mejora del recreo, análisis de uso del laboratorio, evaluación de la cafetería, estudio del tiempo de pantalla del grupo) con la estructura completa y 3-5 páginas};
\end{tikzpicture}
\mbox{}
\newpage

\section*{Apertura: Estructura del informe técnico --- introducción, desarrollo, conclusiones}

\begin{softbox}{milcGradeDark}{milcGradeSoft}{Saber ancestral}
En cualquier cocina del Valle del Cauca, en las casas de los abuelos del Quindío y en las cocinas comunitarias del Pacífico, hay una práctica ancestral que cualquier receta seria respeta: \textbf{toda preparación tiene 3 tiempos}. (1) \textbf{Alistamiento}: antes de prender el fogón, la cocinera saca los ingredientes, los mide, los lava, los pica, los organiza por orden de uso. Sin alistamiento, el aceite se calienta sin la cebolla lista, la sal se olvida, el plato sale apurado. (2) \textbf{Cocción}: aquí pasa lo importante. Se sofríe, se hierve, se condimenta, se prueba, se ajusta. Es la zona más larga del proceso. (3) \textbf{Servido}: la cocinera prueba una última vez, ajusta sal o picante si hace falta, decora si es plato de visita, sirve. Recién entonces el plato existe para el comensal. Quien intenta cocinar saltándose el alistamiento o sin servir bien al final pierde el sabor de todo el trabajo. La receta ancestral exige los 3 tiempos en orden. La sabiduría es simple: \textbf{principio que prepara, medio que cocina, fin que entrega}. El informe técnico moderno hereda esa misma estructura: introducción que prepara, desarrollo que cocina, conclusiones que entregan.
\end{softbox}

\begin{softbox}{milcTurquesa}{milcGris}{Saber contemporáneo}
La \textbf{estructura del informe técnico} es la herramienta más estable del oficio profesional escrito: existe desde los primeros informes científicos del siglo XVII y sigue vigente. Sus 3 secciones obligatorias son: (1) \textbf{Introducción} (15-20 por ciento del informe): presenta el contexto (qué se estudió, por qué importa), el objetivo (qué se pretende mostrar o decidir) y el alcance (qué se cubre y qué queda fuera). Es la sección que orienta al lector apurado. (2) \textbf{Desarrollo} (60-70 por ciento): es el cuerpo del informe. Aquí van la metodología (cómo se hizo el análisis), los datos (qué se encontró), los argumentos (cómo se interpretan los datos). En informes con investigación se subdivide en secciones temáticas con subtítulos. (3) \textbf{Conclusiones} (15-20 por ciento): los hallazgos principales en 1-3 frases claras + recomendaciones accionables (qué se propone hacer con los hallazgos). La regla profesional: \textbf{cada sección tiene función única, no se debe mezclar}. Si pones datos en la introducción o conclusiones sin recomendaciones, el informe pierde rigor. La IA generativa puede ayudar a estructurar y revisar, pero el editor humano decide qué entra en cada sección.
\end{softbox}

\begin{softbox}{milcMostaza}{milcGris}{Pregunta-puente}
\textit{¿Qué sabía la cocinera del Pacífico al respetar los 3 tiempos de la receta, que el novato olvida cuando escribe un informe sin estructura clara? ¿Y por qué un informe con introducción débil pierde al lector aunque el desarrollo sea brillante?}
\end{softbox}

\begin{softbox}{milcVerde}{milcGris}{Saber y saber hacer}
Al terminar podrás: (1) \textbf{identificar} qué información va en cada una de las 3 secciones obligatorias del informe técnico; (2) \textbf{explicar} con tus palabras la función de cada sección y qué pasa cuando se omite o se mezcla; (3) \textbf{aplicar} la estructura completa a un caso real del colegio en informe de 3-5 páginas; (4) \textbf{evaluar} tu informe con criterio de claridad y separación de funciones.
\end{softbox}



\small
\begin{tabularx}{\textwidth}{>{\bfseries}p{3.0cm}Y}
\rowcolor{milcGradePrimary!12}Autor & Dr. Álvaro Cárdenas Orozco\\
DBA & Comunicación profesional con asistencia de IA y herramientas ofimáticas gratuitas (MEN, T\&I)\\
Referentes & Markdown como estándar abierto · Google Workspace · Estoicismo (claridad)\\
Producto final & Plantilla personal reutilizable de informe técnico con las 3 secciones obligatorias estructuradas (introducción con contexto, objetivo y alcance; desarrollo con datos, argumentos y análisis; conclusiones con hallazgos y recomendaciones accionables) + 1 informe técnico aplicado a un caso real del colegio (propuesta de mejora del recreo, análisis de uso del laboratorio, evaluación de la cafetería, estudio del tiempo de pantalla del grupo) con la estructura completa y 3-5 páginas\\
\end{tabularx}
\normalsize

\puente{Antes de empezar, mira el viaje completo. En la sesión 1 distinguiste los 4 tipos de texto profesional. Hoy implementas el \textbf{primero de los 4}: el informe técnico. Las próximas 8 sesiones (Google Docs, Markdown, encuestas, análisis, comercial, correo, mini-proyecto, sustentación) usarán informes técnicos en algún momento. Sin estructura clara hoy, todo lo demás será borroso.}

\begin{titlebox}{milcGradeDark}
Ruta MILC: así vamos a aprender
\end{titlebox}
\begin{tabularx}{\textwidth}{>{\centering\arraybackslash}X>{\centering\arraybackslash}X>{\centering\arraybackslash}X>{\centering\arraybackslash}X}
\phasecard{1}{milcVerde}{milcGris}{Escucha\\[-1mm]\small Observo, escucho y pregunto.} &
\phasecard{2}{milcTurquesa}{milcGris}{Sistematización\\[-1mm]\small Pensamiento computacional.} &
\phasecard{3}{milcMagenta}{milcGris}{Praxis\\[-1mm]\small Construyo y aplico.} &
\phasecard{4}{milcVino}{milcGris}{Evaluación\\[-1mm]\small Reflexión liberadora.}
\end{tabularx}


\puente{Antes de teorizar, vas a hacer una \emph{relectura crítica}: toma el informe técnico real que encontraste en la sesión 1 (Actividad 1) y márcalo con 3 colores distintos: verde para introducción, amarillo para desarrollo, azul para conclusiones. ¿Las 3 secciones están claras? ¿Hay confusiones?}

\begin{titlebox}{milcVerde}
Fase 1 · Escucha
\end{titlebox}

\begin{softbox}{milcVerde}{milcGris}{Escena para escuchar}
\textbf{Actividad 1 · IDENTIFICA --- Relectura crítica del informe real} (15 min · individual). Toma el informe técnico real que encontraste en la sesión 1 (Actividad 1). Imprímelo o ábrelo digital y márcalo con 3 colores distintos: \textbf{verde} para introducción, \textbf{amarillo} para desarrollo, \textbf{azul} para conclusiones. Para cada sección, anota: (1) ¿está claramente identificada (con subtítulo o transición)? (2) ¿qué proporción del informe ocupa? (3) ¿hay mezcla de contenido entre secciones?. Esta lectura crítica te muestra la estructura en uso real.
\end{softbox}

\checkbox Marqué un informe real con 3 colores.\\
\checkbox Identifiqué las 3 secciones.\\
\checkbox Detecté si hay mezcla o claridad.

\begin{infoband}{milcVerde}


\textbf{Lo que va al cuaderno (Actividad 1 · IDENTIFICA).} Título: «Actividad 1 --- 3 secciones bajo la lupa». \textbf{Formato:} captura del informe marcado + observaciones sobre claridad y proporción de cada sección. \textbf{Sabes que terminaste cuando} reconoces la estructura por colores y proporciones.
\end{infoband}

\puente{Ya analizaste 1 informe real. Ahora vas a entender la \textbf{anatomía profesional de las 3 secciones}: qué va en cada una, en qué proporción, los 5 errores típicos del informe novato.}

\begin{titlebox}{milcTurquesa}
Fase 2 · Sistematización con pensamiento computacional
\end{titlebox}

\begin{softbox}{milcTurquesa}{milcGris}{Cuatro pilares aplicados al tema}
Tu lectura crítica te da contexto. Ahora necesitas la \textbf{anatomía profesional} de las 3 secciones del informe técnico.
\end{softbox}

\small
\begin{tabularx}{\textwidth}{>{\bfseries\RaggedRight\arraybackslash}p{3.6cm}Y}
\rowcolor{milcTurquesa!90}\color{white}Pilar & \color{white}\bfseries Aplicación al tema\\
1. Descomposición & Las \textbf{3 secciones obligatorias} se descomponen así: (1) \textbf{Introducción} (15-20 por ciento, ~0.5-1 página en informe de 3-5 pp): \emph{Contexto} (qué problema o tema se aborda, por qué importa ahora), \emph{Objetivo} (qué se pretende mostrar, decidir, recomendar), \emph{Alcance} (qué cubre, qué deja fuera explícitamente). Función: orientar al lector apurado que necesita decidir si seguir leyendo. (2) \textbf{Desarrollo} (60-70 por ciento, ~2-3 páginas): \emph{Metodología} (cómo se obtuvieron los datos: encuesta, observación, registro), \emph{Datos o evidencia} (números, observaciones, citas, gráficos), \emph{Análisis} (interpretación de los datos con argumentos). Puede subdividirse en 2-4 subsecciones según el tema. Función: sostener las conclusiones con evidencia verificable. (3) \textbf{Conclusiones} (15-20 por ciento, ~0.5-1 página): \emph{Hallazgos principales} (en 2-4 frases claras lo más importante encontrado), \emph{Recomendaciones accionables} (qué se propone hacer, quién es responsable, en qué plazo). Función: cerrar con propuesta clara para tomador de decisiones.\\
2. Reconocimiento de patrones & El patrón profesional clave: \textbf{cada sección responde una pregunta distinta}. Introducción responde \emph{``¿de qué trata y por qué importa?''}. Desarrollo responde \emph{``¿qué se encontró y por qué se cree?''}. Conclusiones responden \emph{``¿qué se propone hacer?''}. Si una sección no responde a su pregunta, está mal escrita. La phronesis del oficio consiste en \textbf{no mezclar las preguntas entre secciones}.\\
3. Abstracción & Lo esencial cabe en una regla simple: \textbf{no datos sin análisis, no recomendaciones sin datos}. Si en desarrollo pones solo cifras sin interpretarlas, el lector se pierde. Si en conclusiones recomiendas algo que no está sostenido por los datos del desarrollo, pierdes credibilidad. La estructura es cadena: introducción justifica el estudio, desarrollo aporta evidencia, conclusiones proponen acción con base en esa evidencia.\\
4. Algoritmo conceptual & Algoritmo: (1) construir plantilla personal con las 3 secciones y subsecciones típicas. (2) elegir caso real del colegio. (3) recolectar datos básicos (observación, mini-encuesta, conteo manual). (4) escribir introducción con contexto/objetivo/alcance. (5) escribir desarrollo con metodología/datos/análisis. (6) escribir conclusiones con hallazgos/recomendaciones. (7) verificar proporciones (15/70/15 aproximado). (8) revisar coherencia: ¿las conclusiones se sostienen en los datos del desarrollo?\\
\end{tabularx}
\normalsize

\begin{drawbox}{milcTurquesa}{Anatomía del informe técnico}
\textbf{Introducción (~15-20 por ciento):} contexto + objetivo + alcance. \\[1mm] \textbf{Desarrollo (~60-70 por ciento):} metodología + datos + análisis. \\[1mm] \textbf{Conclusiones (~15-20 por ciento):} hallazgos + recomendaciones. \\[1mm] \textbf{Cadena:} introducción justifica, desarrollo aporta, conclusiones proponen.
\end{drawbox}

\begin{softbox}{milcMostaza}{milcGris}{Errores comunes y su corrección}
(1) \textbf{Introducción muy larga o muy corta}: una página de intro en informe de 3 pp desproporciona; medio párrafo de intro pierde al lector. (2) \textbf{Datos en introducción}: la intro contextualiza, no aporta datos. (3) \textbf{Conclusiones sin recomendaciones}: hallazgos sueltos sin propuesta de acción. (4) \textbf{Recomendaciones sin sustento en datos}: opiniones disfrazadas de conclusión. (5) \textbf{Mezcla de funciones}: análisis en conclusiones, recomendaciones en desarrollo, datos en introducción. (6) \textbf{Sin alcance declarado}: el lector no sabe qué cubre y qué no.
\end{softbox}

\begin{infoband}{milcTurquesa}


\textbf{Lo que va al cuaderno (Actividad 2 · EXPLICA).} Título: «Actividad 2 --- Anatomía del informe técnico». \textbf{Formato:} (a) las 3 secciones con su pregunta clave y proporción; (b) ejemplo de cadena intro-desarrollo-conclusiones; (c) los 6 errores con marca del que más temes. \textbf{Sabes que terminaste cuando} podrías auditar la estructura de cualquier informe real.
\end{infoband}

\puente{Tienes el método. Toca aplicarlo: construyes tu plantilla personal, eliges caso real del colegio, escribes informe técnico completo de 3-5 páginas siguiendo la estructura.}

\begin{titlebox}{milcMagenta}
Fase 3 · Praxis con pensamiento computacional
\end{titlebox}

\begin{softbox}{milcMagenta}{milcGris}{Construyendo el producto con los cuatro pilares}
\textbf{Actividad 3 · APLICA + EVALÚA --- Plantilla personal + informe técnico aplicado} (40 min · individual). Hora de aplicar. Construyes plantilla reutilizable, eliges caso real del colegio, recolectas datos básicos, escribes informe técnico completo de 3-5 páginas.
\end{softbox}

\small
\begin{tabularx}{\textwidth}{>{\bfseries\RaggedRight\arraybackslash}p{3.6cm}Y}
\rowcolor{milcMagenta!90}\color{white}Pilar & \color{white}\bfseries Aplicación al producto\\
Descomposición & Tu plantilla personal tiene 5 elementos: (1) Encabezado del informe (título + autor + fecha + destinatario). (2) Sección Introducción con 3 subsecciones (Contexto, Objetivo, Alcance). (3) Sección Desarrollo con 3 subsecciones (Metodología, Datos, Análisis). (4) Sección Conclusiones con 2 subsecciones (Hallazgos, Recomendaciones). (5) Espacio para firma y fecha al final. Esta plantilla la guardas en Google Docs como plantilla reutilizable.\\
Patrones & Tu informe técnico aplicado se construye en 6 pasos: (1) elegir caso real del colegio (mejora del recreo, uso del laboratorio, estudio del tiempo de pantalla del grupo, evaluación de la cafetería, análisis de inasistencias). (2) recolectar datos básicos esta semana (mini-encuesta a 10-20 personas, conteo manual durante 3 días, observación estructurada). (3) escribir introducción con los 3 elementos. (4) escribir desarrollo con los 3 elementos. (5) escribir conclusiones con los 2 elementos. (6) revisar proporciones, coherencia y la cadena lógica.\\
Abstracción & Sigue el patrón \textbf{``datos verificables sobre opinión''}: en el desarrollo, cada afirmación debe sustentarse en un dato o una observación concreta. \emph{``El recreo es muy corto''} no es dato; \emph{``En encuesta a 15 estudiantes, 12 reportaron que 15 minutos no alcanzan para alimentarse''} sí es dato. La phronesis del informe técnico consiste en \textbf{no afirmar lo que no se puede sostener con evidencia}.\\
Algoritmo de armado & Algoritmo: (1) plantilla. (2) caso. (3) recolección de datos. (4) intro. (5) desarrollo. (6) conclusiones. (7) revisión final: ¿las 3 secciones están proporcionadas?, ¿hay datos que sustentan?, ¿las recomendaciones son accionables?\\
\end{tabularx}
\normalsize

\begin{drawbox}{milcMagenta}{Checklist antes de entregar}
\checkbox Plantilla personal construida y guardada. \\[1mm] \checkbox Caso real del colegio elegido. \\[1mm] \checkbox Datos básicos recolectados. \\[1mm] \checkbox Introducción con contexto + objetivo + alcance. \\[1mm] \checkbox Desarrollo con metodología + datos + análisis. \\[1mm] \checkbox Conclusiones con hallazgos + recomendaciones. \\[1mm] \checkbox Proporciones 15-70-15 respetadas. \\[1mm] \checkbox Cadena lógica verificada.
\end{drawbox}

\begin{softbox}{milcMagenta}{milcGris}{Plantilla de trabajo}
\textbf{Encabezado.} \textit{Título + autor + fecha + destinatario.} \\[1mm] \textbf{Introducción.} \textit{Contexto + Objetivo + Alcance.} \\[1mm] \textbf{Desarrollo.} \textit{Metodología + Datos + Análisis.} \\[1mm] \textbf{Conclusiones.} \textit{Hallazgos + Recomendaciones accionables.} \\[1mm] \textbf{Firma y fecha.} \textit{Cierre formal.}
\end{softbox}

\begin{infoband}{milcMagenta}


\textbf{Lo que va al cuaderno (Actividad 3 · APLICA + EVALÚA).} Título: «Actividad 3 --- Mi informe técnico». \textbf{Formato:} (a) plantilla personal; (b) informe completo de 3-5 páginas con la estructura aplicada; (c) revisión de proporciones y coherencia. \textbf{Sabes que terminaste cuando} tu informe podría ser leído por la coordinación y servir para una decisión.
\end{infoband}

\puente{Lo que sigue resume \emph{qué} entregas hoy: plantilla + informe aplicado. El criterio principal: que cualquier docente o coordinador leyendo tu informe pueda tomar decisión con base en él.}

\begin{titlebox}{milcMostaza}
Producto de la sesión
\end{titlebox}
\textbf{Plantilla personal + informe técnico aplicado de 3-5 páginas}.
\begin{softbox}{milcMostaza}{milcGris}{Criterios de calidad}
(1) Plantilla con las 3 secciones y subsecciones. (2) Caso real del colegio elegido. (3) Datos básicos recolectados. (4) Introducción/desarrollo/conclusiones con proporciones correctas. (5) Recomendaciones accionables. (6) Cadena lógica verificable.
\end{softbox}

\puente{Antes de cerrar, mira la estructura desde las cinco dimensiones humanas. Estructurar antes de escribir es respeto por el lector tomador de decisiones; entregar texto desordenado es traicionar el oficio profesional.}

\begin{titlebox}{milcVino}
Fase 4 · Evaluación liberadora — cinco dimensiones
\end{titlebox}

\begin{softbox}{milcVino}{milcGris}{Sentido de la evaluación}
Evaluar no es solo revisar una nota. Es reconocer qué aprendí, qué cambió en mí, qué emociones manejé, cómo me relaciono con la ciudad y la comunidad, y qué le dejaré a quien viene detrás.
\end{softbox}

\textbf{Mi reflexión liberadora en cinco dimensiones.} Completa cada frase iniciada:

\small
\begin{tabularx}{\textwidth}{>{\bfseries\RaggedRight\arraybackslash}p{3.4cm}Y>{\RaggedRight\arraybackslash}p{4.6cm}}
\rowcolor{milcVino}\color{white}Dimensión & \color{white}\bfseries Mi respuesta (frame starter) & \color{white}\bfseries Algo que descubrí\\
1. Personal & Lo que más cambió en mí fue \linead{6.5cm} & \linead{4.2cm}\\
2. Emocional & La emoción más fuerte fue \linead{2.5cm} y la regulé \linead{3cm} & \linead{4.2cm}\\
3. Ciudadana & En democracia esto importa porque \linead{6cm} & \linead{4.2cm}\\
4. Local (Cartago) & En mi barrio sería útil para \linead{6.5cm} & \linead{4.2cm}\\
5. Intergeneracional & A mi abuela/abuelo le diría \linead{6.5cm} & \linead{4.2cm}\\
\end{tabularx}
\normalsize

\begin{infoband}{milcVino}
\textbf{Para la próxima sustentación.} Vuelve a esta tabla en seis meses. Verás que las respuestas cambian — no porque lo aprendido cambie, sino porque tú cambias. La evaluación liberadora es una conversación contigo a través del tiempo.
\end{infoband}


\puente{Y para terminar, tres voces sobre el orden. Dussel sobre los informes que dignifican al tomador de decisiones, Marco Aurelio sobre la disciplina de la estructura, Floridi sobre el informe como pieza de información cuidada.}

\begin{titlebox}{milcGradeDark}
Triángulo de pensamiento · cierre filosófico
\end{titlebox}

\begin{softbox}{milcVino}{milcGris}{Dussel · lente del nosotros}
\textit{``Un informe técnico estructurado dignifica al tomador de decisiones; uno caótico lo abruma y le quita el tiempo.''}

Quien recibe tu informe (coordinador, rector, gerente futuro) suele tener poco tiempo. La filosofía de la liberación pide ese cuidado: \textbf{estructurar para que el lector tomador de decisiones encuentre lo que necesita rápido}. Las 3 secciones obligatorias permiten que el lector apurado lea solo introducción y conclusiones si no tiene tiempo. La phronesis del informe profesional consiste en \textbf{servir al tomador de decisiones, no abrumarlo}.

\textbf{¿Mi estructura sirve al lector tomador de decisiones, o le exige tiempo que no tiene?}
\end{softbox}
\linead{\linewidth}\\[1mm]
\linead{\linewidth}

\begin{softbox}{milcMostaza}{milcGris}{Estoicismo · Marco Aurelio · cuidado interior}
\textit{``La disciplina de la estructura es virtud del informe; la improvisación es vicio que se nota en la confusión final.''}

Es tentador escribir el informe de corrido sin pensar las secciones. La sabiduría estoica recomienda lo opuesto: \emph{respeta la estructura aunque parezca lenta}. Esa disciplina parece formalismo pero es lo único que produce informes utilizables. Sin estructura el lector se pierde aunque el contenido sea bueno.

\textbf{¿Estoy respetando la estructura con disciplina, o improvisando porque parece más rápido?}
\end{softbox}
\linead{\linewidth}\\[1mm]
\linead{\linewidth}

\begin{softbox}{milcTurquesa}{milcGris}{Floridi · lente de la infoesfera}
\textit{``El informe técnico bien estructurado es pieza de información cuidada en la era del exceso de datos.''}

En la era de la información, los tomadores de decisiones reciben muchos informes. La ética profesional del oficio escrito pide informes bien estructurados que respeten el tiempo del lector. Tu informe técnico de hoy te entrena en una práctica que rige toda profesión técnica: \textbf{producir información cuidada con estructura clara para sostener decisiones reales}.

\textbf{¿Mi informe contribuye a la información cuidada, o agrega ruido al exceso ya existente?}
\end{softbox}
\linead{\linewidth}\\[1mm]
\linead{\linewidth}

\begin{titlebox}{milcVino}
Compromiso final
\end{titlebox}

\small
\begin{tabularx}{\textwidth}{>{\RaggedRight\arraybackslash}p{3.0cm}YYY}
\rowcolor{milcVino}\color{white}\bfseries Criterio & \color{white}\bfseries Lo logré & \color{white}\bfseries En proceso & \color{white}\bfseries Necesito apoyo\\
Comprensión del tema & Explico ideas centrales con mis palabras. & Reconozco ideas pero falta precisión. & Necesito repasar conceptos base.\\
Pensamiento computacional & Apliqué los 4 pilares al diseñar mi producto. & Apliqué algunos pilares. & Necesito releer la fase 2.\\
Aplicación y producto & Mi producto cumple los criterios y se puede revisar. & Producto funcional con ajustes pendientes. & Aún en construcción.\\
Reflexión 5 dimensiones & Respondí cada dimensión con honestidad. & Respondí algunas con detalle. & Mis respuestas son muy generales.\\
Triángulo de pensamiento & Conecté las 3 voces con mi trabajo real. & Conecté al menos una. & No relaciono las voces con mi proceso.\\
\end{tabularx}
\normalsize

\textbf{Hoy entendí que:}\\[1mm]
\linead{\linewidth}\\[1mm]
\linead{\linewidth}

\textbf{Mi compromiso para la próxima sesión es:}\\[1mm]
\linead{\linewidth}\\[1mm]
\linead{\linewidth}

\begin{softbox}{milcMostaza}{milcGris}{Cita de despedida · Marco Aurelio}
\textit{``El obstáculo en el camino se vuelve el camino. Lo que estorba al camino se vuelve camino.''}\\[1mm]
Cada error de hoy, bien observado, se convierte en pieza del camino que pisarás mañana.
\end{softbox}

\small
\begin{tabularx}{\textwidth}{>{\bfseries}p{3.6cm}Y}
\rowcolor{milcGradeDark!12}Nombre completo: & \linead{10cm}\\
Fecha: & \linead{4cm} \quad Curso/grupo: \linead{4cm}\\
Mi firma: & \linead{9cm}\\
\end{tabularx}
\normalsize

\begin{infoband}{milcGradeDark}
\textbf{Para la próxima sesión.} Llega con tu plantilla guardada y el informe técnico aplicado. En la sesión 3 vas a aprender a usar \textbf{Google Docs profesionalmente} con estilos, índice automático y citas. La estructura de hoy es el esqueleto; las herramientas de Docs de mañana son la forma profesional final.
\end{infoband}

\end{document}
